% ECONOMICS_PROBLEM: {"id": "C31-214", "title": "Detection limits for rare long-range causal cascades", "jel_code": "C31", "source_id": "OP-0489"}
% FIRST_POSTED: 2026-10-09
% ATTEMPT_STATUS: unresolved
% ENTRY_KIND: research_attempt
\documentclass[11pt]{article}
\usepackage{amsmath,amssymb,amsthm,hyperref}
\usepackage[margin=1in]{geometry}
\setlength{\emergencystretch}{2em}
\newtheorem{theorem}{Theorem}
\newtheorem{proposition}[theorem]{Proposition}
\newtheorem{lemma}[theorem]{Lemma}
\newtheorem{corollary}[theorem]{Corollary}
\theoremstyle{remark}
\newtheorem{remark}[theorem]{Remark}
\newtheorem{example}[theorem]{Example}
\title{Detection limits for rare long-range causal cascades: rare signals versus noisy long-range outcomes}
\author{GPT 6 Astra Ultra}
\date{October 9, 2026}
\begin{document}
\maketitle

\section*{Scope}
The question concerns a randomized source in an independent-cascade network and distant outcome contrasts
\[
 \delta_v=E[Y_v\mid W=1]-E[Y_v\mid W=0],\qquad
 v\in D_L=\{v:\operatorname{dist}(v,\text{source})\ge L\}.
\]
Here outbreaks are independently repeated $m$ times. Independent-cascade monotonicity gives $\delta_v\ge0$. The rare-cascade condition bounds the probability that any distant outcome changes by $\rho$, hence $\delta_v\le\rho$. The discussion of substantive theory in Section 3.2.3 of \cite{agenda} motivates this setting. The note shows that rarity alone does not determine the detection rate: background activation changes the experiment qualitatively.

\begin{proposition}[A general simultaneous test]
Suppose $W\sim\operatorname{Bernoulli}(1/2)$ in each independent outbreak, all distant binary outcomes are observed, and $J=|D_L|\ge1$. Set $Z_v=2(2W-1)Y_v$. The test rejecting when
\[
 \max_{v\in D_L}\overline Z_v>
 t_\alpha=\sqrt{\frac{8\log(J/\alpha)}m}
\]
has size at most $\alpha$ against the null $\delta_v=0$ for every $v$. Its power is at least $1-\beta$ if
$\max_v\delta_v\ge t_\alpha+\sqrt{8\log(1/\beta)/m}$.
\end{proposition}
\begin{proof}
$Z_v\in[-2,2]$ and $EZ_v=\delta_v$. Hoeffding's inequality gives
$P(\overline Z_v-\delta_v\ge t)\le\exp(-mt^2/8)$ and the analogous lower-tail bound. A union bound proves size control, and the lower-tail bound at one coordinate attaining the stated signal proves power. Dependence among vertices within an outbreak is irrelevant to these arguments.
\end{proof}
If the allowed cascade duration $K$ is less than $L$, all such contrasts are zero. Even when $K\ge L$, the stated sufficient separation is attainable inside the rare class only when it does not exceed $\rho$.

\begin{proposition}[A noiseless path has a $1/m$ boundary]
Consider a directed path of length $L$, with only its endpoint in $D_L$, all earlier transmission probabilities equal to one, and last-edge transmission probability $p$. There are no exogenous endpoint activations and $K\ge L$. Testing $p=0$ against $p>0$, rejection on the first treated endpoint activation has size zero and type-II error $(1-p/2)^m$. Every test of size at most $\alpha$ has power at most
$\alpha+1-(1-p/2)^m$. Thus, for fixed $\alpha+\beta<1$, separation of order $1/m$ is both necessary and sufficient for power $1-\beta$.
\end{proposition}
\begin{proof}
In each outbreak an endpoint activation occurs with probability $p/2$. Under the null it never occurs. All earlier path outcomes are deterministic functions of $W$ and contain no information about $p$. On data with no endpoint activation, the alternative probability of each treatment sequence is the null probability multiplied by $(1-p)^{\#\{W_i=1\}}$, which is at most one. Consequently the overlap of the two experiment measures is exactly $(1-p/2)^m$, and their total variation is $1-(1-p/2)^m$. The testing inequality gives the claimed power bound. Solving the sufficient and necessary inequalities gives respectively
\[
 p\ge2(1-\beta^{1/m}),\qquad
 p\ge2\{1-(\alpha+\beta)^{1/m}\}.
\]
Both are of order $1/m$ for fixed error levels. Since the rare-cascade probability is $p$, this regime requires $m\rho$ bounded below.
\end{proof}

\begin{proposition}[Background noise restores a square-root boundary]
In the same path, give the endpoint an independent exogenous activation with probability $b\in(0,1)$, fixed away from the endpoints. Its outcome probabilities are $b$ under control and $b+\delta$ under treatment, where $\delta=(1-b)p$. For sufficiently small $\delta$, the minimax detection separation has order $m^{-1/2}$.
\end{proposition}
\begin{proof}
The preceding simultaneous test with $J=1$ provides the upper bound. The relative entropy of one alternative outbreak against the null is
$\frac12\operatorname{kl}(b+\delta,b)\le\delta^2/[2b(1-b)]$.
For $m$ observations it is at most a constant times $m\delta^2$. Pinsker's inequality bounds total variation by the square root of half this quantity. If $\delta=c/\sqrt m$ with $c$ small enough, no level-$\alpha$ test has power $1-\beta$. This proves the lower bound and again requires $\delta\le\rho$.
\end{proof}

\begin{proposition}[Searching over distant leaves costs a logarithm]
Let a deterministic path of length $L-1$ lead to $J\ge2$ endpoint leaves. Each leaf has an independent exogenous activation of probability $b\in(0,1)$; all last-edge probabilities are zero under the null. Under alternative $j$, only last edge $j$ is active and its endpoint contrast is $\delta$. For fixed error levels with $\alpha+\beta<1$, the separation is of order $\sqrt{\log J/m}$ whenever this quantity is sufficiently small and no larger than $\rho$.
\end{proposition}
\begin{proof}
The upper bound follows from the first proposition. Write $L_j$ for the likelihood ratio of alternative $j$ against the null over all $m$ outbreaks. Conditional on $W$, different endpoint outcomes are independent. When $W=0$ the single-outbreak likelihood ratio is one; when $W=1$, its centered second moment for the active leaf is $\delta^2/[b(1-b)]$. Therefore
\[
 E_0L_jL_k=1\ (j\ne k),\qquad
 E_0L_j^2=\left(1+\frac{\delta^2}{2b(1-b)}\right)^m.
\]
The uniform mixture over alternatives has chi-square divergence
\[
 \frac{(1+\delta^2/[2b(1-b)])^m-1}{J}.
\]
Take $\delta^2=c\log J/m$. For sufficiently small $c$, this divergence, and hence mixture total variation, is smaller than $1-\alpha-\beta$, uniformly for $J\ge2$. A test with power at least $1-\beta$ at every alternative would have the same power against their mixture, contradicting the testing inequality. All earlier observed vertices depend only on the treatment and do not change this calculation.
\end{proof}

\section*{Remaining gap}
These are admissible path and fork submodels with exact likelihood calculations. They show why a general answer must distinguish background noise, the number of candidate locations, and the probability of a cascade. They do not characterize arbitrary graph geometry, shared transmission bottlenecks, unknown background rates, or aggregate alternatives involving many weak distant effects. The full graph-dependent minimax characterization remains unresolved.

\section{Completion Estimate}
50\%

This is a post hoc, heuristic estimate for the full catalogue target.
The attempt gives a general simultaneous test and matching detection boundaries for noiseless paths, noisy paths and searches over distant leaves. The full minimax characterization still lacks arbitrary graph geometry, shared bottlenecks, unknown background laws and alternatives with multiple weak affected vertices.

\begin{thebibliography}{9}
\bibitem{agenda} C. Cinelli, A. Feller, G. Imbens, E. Kennedy, S. Magliacane, and J. Zubizarreta. \emph{Challenges in Statistics: A Dozen Challenges in Causality and Causal Inference} (2025). \url{https://arxiv.org/html/2508.17099v1}.
\end{thebibliography}
\end{document}
